% ======================================================================
% CHAPTER 1: INTRODUCTION AND RESEARCH RATIONALE
% ======================================================================

\chapter{Introduction and Research Rationale}
\label{chap:introduction}

\section{Research Background and Justification}
\label{sec:background}

\subsection{Meteorological Characteristics of East Asian Sand and Dust Storms and Eco-Health Hazards}
\label{subsec:sds_hazards}

Sand and Dust Storms (SDSs) are catastrophic aeolian phenomena characteristic of arid and semi-arid regions worldwide. They are initiated when turbulent aerodynamic forces lift vast quantities of surface soil, sand grains, and mineral dust into the free atmosphere, resulting in dense particulate suspensions and reducing horizontal visibility below $1{,}000\,\text{m}$~\cite{ref2,zhang2004evolution}. In assessments conducted by the World Meteorological Organization (WMO) and the United Nations Environment Programme (UNEP), sandstorms are classified among the most disruptive transboundary environmental hazards on the planet. Globally, total mineral dust aerosol emissions into the troposphere range between $1{,}000$ and $3{,}000\,\text{Tg}$ per year, fundamentally influencing planetary radiative forcing balance, serving as cloud condensation nuclei, modulating marine biochemical cycles, and inflicting severe damage upon human socioeconomic infrastructure.

East Asia constitutes the second-largest dust-emitting source area and transboundary aerosol pathway globally, surpassed only by the North African Sahara~\cite{ref2,ref3}. The Taklimakan Desert, the Badain Jaran Desert, the Tengger Desert in northwestern China, and the arid Gobi landscapes of southern Mongolia constitute the four primary source cradles of East Asian dust activity. Annually during boreal spring (March to May), the convergence of intense Mongolian cyclogenesis and the southward migration of the Asian Polar Front Jet creates extreme baroclinic shear. Powerful thermal updrafts in the cyclone warm sector, combined with downward momentum transport behind cold fronts, violently dislodge dry fine particulates from denuded surfaces, lifting dust plumes to altitudes of $3{,}000\text{--}8{,}000\,\text{m}$. Under the steering of the mid-tropospheric westerly jet, these dense dust clouds sweep across the Hexi Corridor, the North China Plain, and the Yangtze River Delta, frequently continuing eastward across the Korean Peninsula, the Japanese Archipelago, and into the western Pacific Ocean basin~\cite{ref10,ref16}.

The frequent eruption of East Asian sandstorms triggers extensive multi-dimensional hazards:
\begin{enumerate}
    \item \textbf{Acute Threats to Public Health and Air Quality Degradation}: Suspended dust plumes carry astronomical concentrations of inhalable coarse particulate matter ($\text{PM}_{10}$) and fine particulate matter ($\text{PM}_{2.5}$). Mineral particles are heavily coated with crustal elements, crystalline free silica, and adsorbed biological aerosols, including pathogenic bacteria and fungal spores. For example, during the historic March 15, 2021 super-storm—the most severe event in East Asia in over a decade—ground environmental monitoring networks across Beijing recorded hourly surface $\text{PM}_{10}$ mass concentrations exceeding $8{,}000\,\mu\text{g/m}^3$, with optical visibility collapsing below $500\,\text{m}$ across twelve provinces~\cite{ref2}. Epidemiological studies confirm that severe dust exposure produces acute surges in emergency hospital admissions for respiratory distress (bronchial asthma, chronic obstructive pulmonary disease) and ischemic cardiovascular events.
    \item \textbf{Paralysis of Critical Urban Infrastructure and Transportation Networks}: Extreme particulate loading and severely reduced visual range force major civil aviation hubs to cancel or divert hundreds of flights, while national expressways are closed to prevent multi-vehicle pileups. High-velocity silica dust particles inflict severe mechanical abrasion on photovoltaic solar modules, reducing power generation efficiency by $40\%\text{--}70\%$. Furthermore, the accumulation of conductive mineral dust on ultra-high-voltage transmission insulators induces catastrophic electrical flashover outages under damp, morning mist conditions.
    \item \textbf{Staggering National Economic Losses}: According to emergency management assessments, annual direct and indirect economic losses attributable to sandstorm disasters across northern China range between $10$ and $15\,\text{billion RMB}$ ($\sim \$1.5\text{--}2.2\,\text{billion USD}$)~\cite{ref2,ref7}. Prolonged topsoil erosion depletes agricultural soil nutrients, destroys protective mulch and greenhouse facilities, and accelerates secondary desertification along fragile ecological ecotones.
\end{enumerate}

\subsection{Alignment with National Ecological Security and Disaster Mitigation Strategies}
\label{subsec:national_strategy}

Developing advanced, long-lead-time sandstorm forecasting capabilities aligns directly with critical national strategic initiatives:
\begin{enumerate}
    \item \textbf{Scientific Grounding for the National ``Three-North'' Shelterbelt Program}: The ongoing sixth phase of the ``Three-North'' Shelter Forest Program represents a monumental environmental shield spanning China's northwestern, northern, and northeastern desert margins. Accurate medium- to long-term projections of soil moisture deficits, aerodynamic wind regimes, and dust saltation pathways are indispensable for optimizing the spatial arrangement of sand-fixing vegetation barriers and quantitatively assessing long-term ecological rehabilitation returns.
    \item \textbf{Paradigm Shift in Comprehensive Disaster Prevention and Reduction}: In accordance with the National Comprehensive Disaster Prevention and Reduction Plan, disaster risk management is undergoing a structural paradigm shift from ``passive post-disaster relief'' toward ``proactive pre-disaster risk mitigation.'' Providing accurate predictions 3 to 15 days in advance equips municipal administrations, environmental sanitation agencies (e.g., targeted water misting for airborne dust suppression), and transportation authorities with multi-day preparation buffers, dramatically reducing economic exposure and safeguarding human lives.
    \item \textbf{Transboundary Ecological Security for the Green Belt and Road Initiative}: Dust storms are inherently transboundary meteorological phenomena. The border regions linking southern Mongolia with Inner Mongolia represent major transboundary emission corridors. High-precision intelligent forecasting provides actionable scientific support for collaborative international environmental stewardship within the China-Mongolia Desertification Prevention and Control Center, safeguarding regional ecological stability along the Silk Road economic belt.
\end{enumerate}

\subsection{Operational Bottlenecks in Extended-Range (3--15 Day) Forecasting}
\label{subsec:extended_range_bottleneck}

While modern meteorological operations deliver dependable deterministic forecasts within short-range horizons (1 to 3 days, or 24 to 72 hours), operational forecasting skill experiences a precipitous and catastrophic decline when moving into the **extended-range window (3 to 15 days)**~\cite{ref2,ref10}:
\begin{enumerate}
    \item \textbf{Sensitive Dependence on Initial Conditions and the Chaotic Lyapunov Horizon}: Atmospheric fluid dynamics is non-linear and chaotic. According to Lorenz's foundational turbulence theory~\cite{lorenz1963deterministic}, infinitesimal perturbations in initial atmospheric measurements amplify exponentially over phase space. Beyond 72 hours, rapid error propagation degrades raw Numerical Weather Prediction (NWP) wind fields, causing modeled dust advection trajectories to diverge drastically from observational reality.
    \item \textbf{Coarse Topographic Smoothing over Narrow Orographic Corridors}: East Asian dust plumes are channeled through narrow, complex mountain gaps (such as the Hexi Corridor, the Helan Mountain Gorge, and the Yan-Taihang Mountain passes). Global and regional NWP models operating at grid resolutions of $0.1^\circ$ to $0.25^\circ$ mechanically smooth steep topographic features. This spatial flattening fails to resolve sub-grid venturi wind acceleration and low-level jet development, resulting in severe temporal phase lags and systematic underestimation of downstream dust peak arrivals.
    \item \textbf{Failure of Static Empirical Emission Schemes During Spring Thaw}: Conventional chemical transport models parameterize surface dust fluxes using empirical formulations (e.g., Gillette, Shao, Kok) calibrated on idealized wind-tunnel observations~\cite{gillette1979environmental,shao2008physics,kok2012physics}. These schemes assume stationary aerodynamic roughness and fixed soil texture. However, during early spring, rapid soil thawing, snowmelt retreat, and fluctuating topsoil moisture generate dramatic, non-linear shifts in the threshold friction velocity, which static empirical schemes are fundamentally incapable of capturing.
\end{enumerate}

---

\section{State-of-the-Art and Scientific Gaps}
\label{sec:literature_summary}

\subsection{Numerical Weather Prediction (NWP) in Dust Modeling: Capabilities and Shortcomings}
\label{subsec:nwp_limits}

Contemporary operational forecasting systems—including the European Centre for Medium-Range Weather Forecasts Integrated Forecasting System (ECMWF IFS-AER), the China Meteorological Administration Global Forecast System (CMA-GFS), and Eulerian chemical transport models such as WRF-Chem and CUACE/Dust—simulate dust lifecycles by numerically solving planetary momentum conservation and aerosol continuity partial differential equations (PDEs)~\cite{ref2}:
\begin{equation}
\label{eq:aerosol_transport_intro}
\frac{\partial C}{\partial t} + \nabla \cdot (\mathbf{u} C) = \nabla \cdot (\mathbf{K} \nabla C) + S_{\text{emission}} - v_d C - \Lambda_{\text{scav}} P_{\text{rain}} C
\end{equation}
where $C$ denotes particulate mass concentration, $\mathbf{u}$ is the three-dimensional advective wind vector, $\mathbf{K}$ is the turbulent eddy diffusivity tensor, $S_{\text{emission}}$ represents the surface saltation emission flux, $v_d$ is the dry gravitational settling velocity, and $\Lambda_{\text{scav}} P_{\text{rain}}$ accounts for precipitation-induced wet scavenging.

Although dynamical numerical models offer indisputable physical interpretability, their operational execution beyond 72 hours reveals severe structural bottlenecks. Integrating global high-order PDEs requires massive supercomputing clusters, consuming hundreds of thousands of CPU core-hours per ensemble run. This computational burden prevents real-time, rolling-horizon updates during sudden emergency crises. Furthermore, unresolved sub-grid turbulent wind bursts must be parameterized via empirical approximations, causing cumulative numerical errors to compound monotonically with lead time~\cite{ref2,ref3}.

\subsection{Machine Learning and Deep Spatiotemporal Networks in Atmospheric Sciences}
\label{subsec:ml_weather_progress}

Driven by the exponential growth of multi-source earth observation data (satellite remote sensing, sounding profilers, surface telemetry, and reanalysis archives), data-driven machine learning (ML) has emerged as a transformative paradigm across atmospheric and environmental sciences~\cite{ref1,ref32,ref40}.

In sandstorm forecasting, pioneering studies introduced shallow learners, including Multilayer Perceptrons (MLP)~\cite{ref3} and Support Vector Machines (SVM)~\cite{ref4}, to classify dust occurrence at discrete observation stations. To address the acute **class imbalance** stemming from the rarity of severe dust storms relative to calm days, Xie et al.~\cite{ref5} and Zhang et al.~\cite{ref6} developed hybrid adaptive sampling and SMOTE-AdaBoost architectures, markedly improving the recall of rare extreme events. Subsequently, Kaboodvandpour et al.~\cite{ref7} and Shaiba et al.~\cite{ref9} systematically demonstrated the superiority of ensemble tree models (Random Forest, Gradient Boosting) over linear statistical baselines.

Over the past five years, deep spatiotemporal neural networks have achieved remarkable breakthroughs:
\begin{itemize}
    \item Recurrent architectures embedding spatial convolutions (ConvLSTM~\cite{ref17}, PredRNN~\cite{ref18}) and spatiotemporal attention models (Earthformer~\cite{ref20}, Rainformer~\cite{ref19}) have redefined precipitation nowcasting and cloud motion extrapolation;
    \item Breakthrough global foundation models, such as the Aerosol-Meteorology Coupled Large Model Forecasting System (AI-GAMFS), achieved operational 5\,km resolution forecasts with lead times reaching 120 hours (5 days), reducing East Asian dust forecast errors by $38\%\text{--}74\%$ relative to ECMWF IFS~\cite{ref2};
    \item Transfer learning and multi-branch architectures (INB-CNN~\cite{ref10}, Stacking LSTM-CNN~\cite{ref16}) demonstrated substantial performance gains in regional dust trajectory tracking.
\end{itemize}

Nonetheless, standard data-driven networks suffer from a critical pathology: **physical hallucinations**. When trained purely on empirical loss minimization, unconstrained networks memorize spurious statistical correlations in training corpora. Consequently, they frequently predict dust generation in completely calm, low-wind conditions or generate mass out of a vacuum, violating basic aerodynamic thresholding and mass continuity principles~\cite{raissi2019physics}.

\subsection{Physics-Informed Neural Networks (PINNs) in Fluid Mechanics and Environmental Transport}
\label{subsec:pinn_progress}

To bridge the dichotomy between black-box statistical learning and dynamical rigor, Raissi et al.~\cite{raissi2019physics} formulated Physics-Informed Neural Networks (PINNs). By leveraging reverse-mode automatic differentiation (AD), PINNs embed the differential operators governing physical conservation laws directly into the neural network loss function:
\begin{equation}
\label{eq:pinn_loss_intro}
\mathcal{L}_{\text{total}}(\theta) = \mathcal{L}_{\text{data}}(\theta) + \lambda_{\text{phys}} \mathcal{L}_{\text{PDE}}(\theta)
\end{equation}
In canonical fluid mechanics problems (e.g., incompressible Navier-Stokes flows, porous media transport), PINNs demonstrate remarkable robustness and extrapolation fidelity under sparse data regimes. However, within continental-scale atmospheric dust modeling, the mathematical challenge of embedding discontinuous, highly non-linear saltation thresholds alongside 2D advective conservation PDEs within deep spatiotemporal graphs has remained an unaddressed frontier.

\subsection{Geospatial Exposure Analytics and Disaster Behavioral Response}
\label{subsec:socio_progress}

In disaster science, atmospheric engineering has traditionally operated in strict isolation from social science and public emergency management. Even a theoretically flawless meteorological forecast provides negligible societal value if municipal authorities lack operational playbooks and citizens fail to translate forecast warnings into protective compliance.

Geospatially, Tobler's First Law of Geography~\cite{tobler1970cellular} establishes that environmental hazards exhibit strong spatial autocorrelation. Spatial statistical techniques, including Global Moran's $I$ and Local Indicators of Spatial Association (LISA), provide mathematical tools for delineating regional risk hotspots. In behavioral psychology, Lindell and Perry's Protective Action Decision Model (PADM)~\cite{lindell2012protective} formalizes the cognitive-affective causal chain from environmental hazard cues to perceived risk, institutional trust, and protective behavioral compliance. Nevertheless, empirical frameworks that seamlessly link numerical AI prediction uncertainty envelopes ($[P_{10}, P_{90}]$) to GIS exposure mapping and PADM structural equation modeling (SEM) are completely absent in the current literature.

---

\section{Research Objectives and Core Scientific Questions}
\label{sec:objectives_questions}

\subsection{Research Objectives}
\label{subsec:objectives}

To resolve the dual bottlenecks of extended-range numerical skill degradation and unconstrained artificial intelligence hallucinations, this dissertation pursues the following primary research objectives:
\begin{enumerate}
    \item Develop an end-to-end, multi-source data-driven methodology tailored for East Asian sandstorm forecasting across the 3--15 day extended-range horizon;
    \item Establish a dual-line predictive architecture uniting Main Line A (tree-based NWP statistical residual correction with quantile uncertainty) and Main Line B (an end-to-end deep spatiotemporal foundation model coupling AI-GAMFS priors, 14-node corridor graph diffusion, and PINN conservation regularizers);
    \item Formulate a continuous mathematical representation of Owen's aerodynamic saltation threshold and 2D advective mass continuity within the backpropagation computation graph, eliminating non-physical predictions;
    \item Synthesize physical forecasting outputs with GIS geostatistical kriging and PADM structural equation modeling, delivering an operational 4-tier municipal emergency response decision matrix.
\end{enumerate}

\subsection{Core Scientific Questions}
\label{subsec:scientific_questions}

To achieve these objectives, this research addresses three central scientific questions:
\begin{itemize}
    \item \textbf{Scientific Question 1 (Non-linear Topographic Bias Correction):}
    \textit{How can machine learning architectures effectively capture and eliminate the systematic, non-linear spatiotemporal residuals of NWP models over complex mountainous terrain without destabilizing dynamic atmospheric balance?}
    \item \textbf{Scientific Question 2 (Physical Prior Embedding and Conservation Constraints):}
    \textit{How can aerodynamic saltation thresholds and advective mass conservation PDEs be mathematically embedded into deep neural networks to guarantee physical consistency and prevent unconstrained hallucinations during extended extrapolation?}
    \item \textbf{Scientific Question 3 (Spatial Hazard Transmission and Behavioral Compliance):}
    \textit{What are the quantitative transmission pathways through which spatiotemporal dust exposure patterns affect public risk perception, institutional trust, and protective compliance? How can forecast uncertainty intervals be directly coupled with municipal emergency actions?}
\end{itemize}

---

\section{Main Research Contents and Master Technical Roadmap}
\label{sec:contents_roadmap}

\subsection{Four Key Research Modules}
\label{subsec:four_modules}

The dissertation is structured into four integrated, cumulative modules:

\begin{enumerate}
    \item \textbf{Module 1: Multi-Source Heterogeneous Data Harmonization and Feature Engineering (Chapter 3)} \\
    A comprehensive data collection pipeline ingests multi-source data covering East Asia ($73^\circ\text{E}\text{--}135^\circ\text{E}, 18^\circ\text{N}\text{--}54^\circ\text{N}$) from 2018 to 2025. Data streams include NWP forecasts (ECMWF IFS, CMA-GFS), high-resolution reanalysis (ERA5), multi-spectral satellite retrievals (MODIS, FY-4A), and ground telemetry from national ambient air quality monitoring networks. Automated pipelines perform spatio-temporal alignment onto a standardized 5\,km WGS84 grid, compiling a 32-dimensional physics feature tensor encompassing dust emission potential, hydrodynamic transport, boundary layer thermodynamics, and upstream historical memory.

    \item \textbf{Module 2: Dual-Line Machine Learning Predictive Framework and PINN Integration (Chapter 4)} \\
    \begin{itemize}
        \item \textit{Main Line A (NWP Statistical Residual Correction)}: Implements LightGBM, CatBoost, and XGBoost ensembles optimized via cost-sensitive focal loss to mitigate extreme-event class imbalance. Asymmetric pinball quantile loss functions output calibrated prediction bounds ($P_{10}, P_{50}, P_{90}$).
        \item \textit{Main Line B (Deep Spatiotemporal PINN Foundation Model)}: Couples AI-GAMFS foundation representations with a 14-node East Asian transport corridor Spatiotemporal Graph Neural Network (ST-GNN) employing 3-support Chebyshev polynomial diffusion. A custom PINN regularizer embeds Owen's aerodynamic saltation gate and advective mass continuity into the loss function.
        \item \textit{Adaptive Lead-Time Blending}: An adaptive sigmoid decay gate $w(t) = 1 / (1 + e^{\beta (t - \tau)})$ dynamically balances predictions, leveraging NWP dynamic skill at short horizons while smoothly transferring dominance to deep physical representations at extended lead times (3--15 days).
    \end{itemize}

    \item \textbf{Module 3: Geospatial Exposure Modeling and Societal Behavioral SEM (Chapter 5)} \\
    Ordinary Kriging with fitted empirical semivariograms interpolates discrete station telemetry into continuous spatial exposure surfaces. Global Moran's $I$ and LISA statistics identify regional High-High hazard hotspots. Based on the Protective Action Decision Model (PADM), a large-scale empirical survey ($N=842$ validated responses across northern China) is analyzed in SPSS and AMOS, establishing a Structural Equation Model (SEM) that quantifies how environmental hazard cues, mediated by institutional trust, drive protective behavioral compliance.

    \item \textbf{Module 4: Historical Retrospective Validation, Model Explainability, and Operational Playbooks (Chapter 6)} \\
    The integrated framework is validated across benchmark historical events, including the March 2021 mega-storm and the April 2025 Sichuan Basin intrusion. Tree SHAP and graph attention heatmaps illuminate physical driving mechanisms, including snowpack retreat, cyclogenesis, and orographic jet channeling. Finally, a 4-tier municipal emergency decision matrix links forecast uncertainty bounds directly to operational municipal interventions.
\end{enumerate}

\subsection{Master Technical Roadmap}
\label{subsec:technical_workflow}

The overarching technical roadmap of this dissertation is illustrated in Figure~\ref{fig:technical_roadmap}.

\begin{figure}[H]
\centering
\begin{tikzpicture}[
    scale=0.88, transform shape,
    box/.style={rectangle, draw=ustbblue, fill=ustbblue!8, text centered, rounded corners=4pt, minimum height=0.85cm, line width=1pt, font=\small},
    decision/.style={diamond, draw=brickred, fill=brickred!10, text centered, line width=1pt, font=\small, aspect=1.8},
    pillar/.style={rectangle, draw=forestgreen, fill=forestgreen!8, text centered, rounded corners=4pt, line width=1pt, font=\small},
    arrow/.style={-{Stealth[scale=1.1]}, thick, draw=navyblue},
    darrow/.style={thick, dashed, draw=gray}
]

% Level 1: Demand & Objective
\node[box, fill=ustbblue!20, minimum width=13cm] (strat) {\textbf{National Strategic Imperatives: Disaster Mitigation, Three-North Program \& Belt and Road Security}};

% Level 2: Multi-Source Data
\node[box, below=0.6cm of strat, minimum width=13cm] (data) {\textbf{Multi-Source Data Ingestion and Spatiotemporal Alignment (Chapter 3)}\\ \footnotesize • ECMWF / CMA-GFS Numerical Forecasts \quad • ERA5 Atmospheric Reanalysis \quad • MODIS / FY-4A Remote Sensing \quad • 5\,km WGS84 Grid};

% Level 3: Dual Lines
\node[box, below left=1.0cm and 0.2cm of data, text width=6.2cm] (lineA) {\textbf{Main Line A: NWP Residual Learning (Chapter 4)}\\ \footnotesize • LightGBM / CatBoost Gradient Boosting\\ • Cost-Sensitive Focal Loss for Rare Extremes\\ • Asymmetric Quantile Regression ($P_{10}, P_{50}, P_{90}$)};

\node[box, below right=1.0cm and 0.2cm of data, text width=6.2cm] (lineB) {\textbf{Main Line B: Deep Spatiotemporal PINN (Chapter 4)}\\ \footnotesize • AI-GAMFS Aerosol-Meteorology Prior Feature Store\\ • 14-Node East Asian Corridor ST-GNN\\ • 3-Support Chebyshev Directed Graph Diffusion};

% Level 4: PINN Injection & Lead-time Blending
\node[box, fill=brickred!15, draw=brickred, below=2.7cm of data, minimum width=13cm] (pinn) {\textbf{Physics-Informed Regularization (PINN) and Adaptive Lead-Time Blending}\\ \footnotesize • Owen (1964) Aerodynamic Saltation Threshold Loss $\mathcal{L}_{\text{Owen}}$ \quad • 2D Advective Mass Conservation PDE Penalty $\mathcal{L}_{\text{PDE}}$\\ \footnotesize • Dynamic Lead-Time Blending Weight Function $w(t)$ Transitioning from Line A to Line B};

% Level 5: Downstream Branches
\node[pillar, below left=1.0cm and 0.2cm of pinn, text width=6.2cm] (gis_socio) {\textbf{Spatial Exposure \& Behavioral Response (Chapter 5)}\\ \footnotesize • Ordinary Kriging Continuous Spatial Exposure\\ • Global Moran's $I$ \& LISA High-High Hotspots\\ • PADM Structural Equation Modeling ($N=842$)};

\node[pillar, below right=1.0cm and 0.2cm of pinn, text width=6.2cm] (case_xai) {\textbf{Historical Hindcasting \& Operational Playbook (Chapter 6)}\\ \footnotesize • March 2021 Super-Storm 15-Day Extended Hindcast\\ • Tree SHAP Attribution \& Orographic Jet Diagnostics\\ • 4-Tier Municipal Emergency Response Decision Matrix};

% Level 6: Conclusion
\node[box, fill=ustbblue!20, below=2.7cm of pinn, minimum width=13cm] (concl) {\textbf{Conclusions, Methodological Synthesis, and Operational Recommendations (Chapter 7)}};

% Connectors
\draw[arrow] (strat) -- (data);
\draw[arrow] (data.south) -| (lineA.north);
\draw[arrow] (data.south) -| (lineB.north);
\draw[arrow] (lineA.south) |- (pinn.north west);
\draw[arrow] (lineB.south) |- (pinn.north east);
\draw[arrow] (pinn.south) -| (gis_socio.north);
\draw[arrow] (pinn.south) -| (case_xai.north);
\draw[arrow] (gis_socio.south) |- (concl.north west);
\draw[arrow] (case_xai.south) |- (concl.north east);

\end{tikzpicture}
\caption{Master Technical Roadmap of the Dissertation}
\label{fig:technical_roadmap}
\end{figure}

---

\section{Major Research Innovations}
\label{sec:innovations}

This dissertation formulates three explicit theoretical and engineering innovations:

\begin{enumerate}
    \item \textbf{Innovation 1: Physics-Informed Neural Network (PINN) Constraint Mechanism Embedding Aerodynamic Saltation Thresholds and Advective Mass Continuity} \\
    \textit{Eliminates the persistent pathology of unphysical hallucinations (e.g., predicting sandstorms under calm wind conditions) in deep learning models.} \\
    Standard neural networks optimize purely empirical empirical risk, leading to physically impossible predictions during out-of-distribution extrapolation. This research formulates a differentiable mathematical representation of Owen's (1964) saltation threshold criterion ($u_* > u_{*t}$) and a 2D advective mass continuity PDE directly within the PyTorch backpropagation graph. A dynamic multi-objective balancing algorithm (GradNorm) stabilizes gradient magnitudes between empirical loss and higher-order differential residuals. Empirical evaluations demonstrate that this mechanism attains a **$98.2\%$ physical law compliance rate**, reducing root mean square error (RMSE) by 18.2\% in unobserved extreme wind regimes.

    \item \textbf{Innovation 2: Bidirectional Spatiotemporal Graph Neural Network (ST-GNN) Mapping East Asian Dust Transport Corridor Topology} \\
    \textit{Overcomes the spatial distortion, phase lag, and computational burden caused by coarse Euclidean grid smoothing over narrow mountain passes.} \\
    Rather than treating complex topography as a uniform flat 2D grid, this study models the East Asian dust transport corridor as a directed 14-node topological graph spanning source regions, orographic funnels (e.g., Hexi Corridor, Helan Mountain pass), and downstream megalopolises. A bidirectional ST-GNN utilizing 3-support Chebyshev polynomial diffusion (self-connection $P_0$, downstream advection $P_1$, upstream back-tracking $P_2$) explicitly resolves orographic jet acceleration. This architecture achieves an 8-fold reduction in computational training cost compared to full-grid 3D-CNNs, while extending the effective forecasting horizon from 3 days to **10--15 days**.

    \item \textbf{Innovation 3: Tri-Pillar Integrated Framework Bridging Atmospheric AI Prediction, Spatial Risk Exposure, and Societal Behavioral Response} \\
    \textit{Resolves the long-standing disconnect between meteorological engineering forecasts and societal emergency mitigation.} \\
    This research establishes an interdisciplinary paradigm uniting machine learning forecasting, GIS spatial econometrics (Ordinary Kriging, Global Moran's $I$, LISA), and psychometric Structural Equation Modeling (SEM) grounded in the Protective Action Decision Model (PADM). Utilizing an empirically validated survey dataset ($N=842$) across northern China, the model quantifies the transmission pathways between physical hazard cues, institutional trust, and public protective compliance. This provides municipal authorities with an operational decision matrix that maps calibrated forecast uncertainty intervals ($[P_{10}, P_{90}]$) to actionable, four-tier municipal emergency interventions.
\end{enumerate}

---

\section{Thesis Structure and Chapter Organization}
\label{sec:thesis_organization}

The dissertation is structured into seven coherent chapters:

\begin{itemize}
    \item \textbf{Chapter 1: Introduction and Research Rationale}: Establishes the research background, public health and economic hazards of East Asian sandstorms, alignment with national ecological programs, extended-range forecasting bottlenecks, core scientific questions, technical roadmap, and key innovations.
    \item \textbf{Chapter 2: Literature Review and Theoretical Foundations}: Synthesizes atmospheric boundary layer dynamics, Navier-Stokes balance, Owen saltation mechanics, aerosol transport PDEs, NWP limitations, the Lyapunov horizon, machine learning in atmospheric sciences, PINN formulations, spatial econometrics, and PADM theory, delineating the three critical research gaps.
    \item \textbf{Chapter 3: Multi-Source Data Harmonization and Feature Engineering}: Details data acquisition pipelines (ECMWF, CMA-GFS, ERA5, MODIS, FY-4A, station telemetry), 5\,km WGS84 grid alignment, quality control protocols, and the construction of the 32-dimensional physics feature tensor.
    \item \textbf{Chapter 4: Dual-Line Machine Learning Predictive Framework and PINN Architecture}: Formulates Main Line A (tree ensembles with quantile regression), Main Line B (AI-GAMFS priors, 14-node ST-GNN, and PINN conservation loss functions), and the adaptive lead-time blending gate.
    \item \textbf{Chapter 5: Spatial Risk Exposure Mapping and Societal Behavioral Response}: Details Ordinary Kriging continuous surface interpolation, Moran's $I$ and LISA spatial clustering, empirical survey administration ($N=842$), and AMOS Structural Equation Modeling of protective action compliance.
    \item \textbf{Chapter 6: Historical Event Validation, Interpretability, and Operational Playbooks}: Evaluates framework performance against historical benchmark events (March 2021 super-storm and April 2025 basin intrusion), conducts Tree SHAP feature attribution, and delivers the 4-tier municipal emergency response decision matrix.
    \item \textbf{Chapter 7: Conclusions and Outlook}: Synthesizes major scientific findings, summarizes engineering innovations, addresses study limitations, and outlines future research trajectories in atmospheric foundation models and cross-disciplinary disaster mitigation.
\end{itemize}
